%==============================================================================
%  PlantMetWiki -- COVER LETTER
%  Compiles to its own PDF, uploaded as the "covering letter" file.
%  Scientific Data: uploading a covering letter is a technical requirement of
%  the submission system; its content is not used to make editorial decisions.
%
%  SUBMISSION PORTAL:
%  https://submission.springernature.com/submission/ed40d312-e1b7-4d02-9075-86bd6970525f
%
%  TIMELINE
%    16 October 2026 -- co-author approval deadline and final comments
%    26 October 2026 -- submission to the journal
%
%  TWO FIGURES WERE CHANGED FROM THE SUPPLIED DRAFT so the letter agrees with
%  the manuscript and the live endpoint -- see the notes at the end of this file.
%
%  The journal's note on suggested reviewers (not part of the letter):
%    "Suggested reviewers or Editorial Board Members: please include these as
%     you wish, but neither are mandated, and rest assured we have processes to
%     assign manuscripts and reviewers without input, and suggestions are not
%     required to expedite review."
%==============================================================================
\documentclass[11pt,a4paper]{article}

\usepackage{iftex}
\ifPDFTeX
  \usepackage[utf8]{inputenc}
  \usepackage[T1]{fontenc}
  \usepackage{textcomp}
\else
  \usepackage{fontspec}
\fi

\usepackage[margin=2.5cm]{geometry}
\usepackage[hidelinks]{hyperref}
\usepackage{setspace}

\onehalfspacing
\setlength{\parindent}{0pt}
\setlength{\parskip}{0.8em}
\pagestyle{empty}

\begin{document}

\hfill \today

\vspace{1em}
To the Editors,\\
\textit{Scientific Data}

\vspace{0.5em}
\textbf{Re: submission to the Collection ``10 years of the FAIR principles''}\\
\textbf{A FAIR knowledge graph for plant metabolic pathway cross-species
representation and integration}

\vspace{1em}
Dear Editors,

We would like to submit our manuscript for the collection ``10 years of the FAIR principles''.

Mention the preprint 

Plants represent the most chemically diverse kingdom of life thanks to their extremely diverse
specialized metabolism. Knowledge of biosynthesis and its evolution is key to discovering new plant
specialized metabolites with bioactivities relevant for human applications. However, only a few biosynthetic pathways of plant specialized metabolites
are fully resolved, limiting the potential to metabolically engineer these compounds at an
industrial scale for applications in medicine, novel foods, agriculture, and bioenergy. Furthermore,
plant pathway knowledge is fragmented across different databases and individual publications, and
the main reference resource, PlantCyc, is distributed as flat files that cannot be queried together
with other resources. Plant pathway data are therefore findable, but not interoperable or readily
reusable, and we are not aware of an open-access infrastructure that integrates pathway predictions
and extends biochemical reaction knowledge with natural product information.

Here, we address these gaps with Plant Metabolic Pathways Wiki (PlantMetWiki), a FAIR Linked Open
Data knowledge graph of plant biosynthetic pathways from PlantCyc (Hawkins et al., 2025;
doi:10.1093/nar/gkae1231) connecting information to known and predicted genomic loci of clustered
biosynthetic genes from plantiSMASH 2.0 (Del Pup et al., 2026; doi:10.1016/j.jmb.2026.169798) and
MIBiG (Zdouc et al., 2025; doi:10.1093/nar/gkae1115), respectively. PlantMetWiki includes 19,927
unique biochemical conversions across 1,162 curated pathways from 424 plant species or genera, and
leverages a novel approach to model multispecies pathways by adapting the existing WikiPathways
infrastructure (Agrawal et al., 2024; doi:10.1093/nar/gkad960). We do note that a previous version of our manuscript has been posted to bioRxiv - add ref.

In our manuscript we describe how we implemented each of the FAIR principles, including persistent identifiers for all data elements, VoID and PROV-O provenance for every
named graph, the vocabularies we use and define published alongside the data so that the model can
be resolved from the endpoint itself, and versioned Zenodo releases for each step of the pipeline.
We also show interoperability via federated queries to Wikidata (Waagmeester et al., 2020;
doi:10.7554/eLife.52614), which led to identifying 421 metabolites without any structurally related entry. Identifier mapping with BridgeDb materializes cross-references for 3,198 of the 4,577
metabolites (70\%), raising ChEBI coverage from 174 to 2,805 metabolites and PubChem from 75 to
2,926. Using PlantMetWiki we also quantify the potential of translating knowledge across across species, by classifying 29,061
species--reaction combinations into annotated, inferable, and non-inferable annotation gaps.

PlantMetWiki is available in Semantic Web format and can be accessed via a dedicated SPARQL endpoint
for programmatic access, and via the SNORQL user interface with prewritten SPARQL queries to analyze
pathway information in the knowledge graph and combine it with user data, such as (multi-)omics
datasets, phenotypes, and predicted pathways. The results can be downloaded as CSV, JSON, or XML for
further visualization of pathway annotations, hypothesis generation, and natural product discovery.
Our approach enables the generation of
hypotheses from plant metabolism to other research areas, including biomarker discovery,
toxicology, precision nutrition, and forensic (wildlife) sciences.

Our conversion workflow operates on standard BioCyc flat-file exports, so the same steps can be
reused for MetaCyc and other BioCyc databases featuring other kingdoms than plants. We anticipate that PlantMetWiki will support
collaborative curation and hypothesis generation in plant biosynthesis by leveraging linked data,
thereby advancing the discovery and characterization of plant natural products.

Editorial Board Member: 
Philippe Rocca-Serra

Suggested reviewers: 
Robin Haw
Pierre Marie Allard, Affiliation, email 
Pierre Larmande
https://bioinformatics.ca/instructors/robin-haw/
Sebastien Moretti
Disha Tandon


\vspace{1em}
On behalf of the PlantMetWiki coauthors,

\vspace{0.5em}
Elena Del Pup, Max Muller, Marvin Martens, Egon Willighagen, Marnix H. Medema,
Denise Slenter \& Justin J.J. van der Hooft

\end{document}

%==============================================================================
%  CHANGES FROM THE SUPPLIED DRAFT -- revert if you disagree
%
%  1. "over 11,000 unique biochemical reactions from 643 plant species or
%     genera"  ->  "19,927 unique biochemical conversions across 1,162 curated
%     pathways from 424 plant species or genera".
%
%     Checked against the live endpoint:
%       - No measure in the graph is near 643 taxa. The graph holds 424
%         NCBITaxon identifiers (423 species/genus level plus Viridiplantae);
%         PlantCyc 17.0 supplies 439. The manuscript Results state 424.
%         643 would contradict the manuscript.
%       - Nothing labelled "reactions" is near 11,000. The closest figure is
%         11,138 unique DataNodes -- which are biological entities (genes,
%         enzymes, metabolites, complexes), not reactions. Reactions are
%         19,927 wp:Conversion edges, or 1,316 standalone RC* reaction files.
%         "11,000 reactions" looks like the entity count mislabelled.
%
%  2. Citation years aligned with references.bib and the published versions:
%       - plantiSMASH 2.0: the draft cited the 2025 bioRxiv preprint
%         (doi:10.1101/2025.10.28.683968). It is now published in J Mol Biol
%         (2026; doi:10.1016/j.jmb.2026.169798), which is what the manuscript
%         cites. Updated.
%       - MIBiG: Zdouc et al. 2024 -> 2025, matching NAR 53 D678-D690 and the
%         manuscript's reference list. DOI unchanged.
%==============================================================================
